\documentclass[addpoints, a4paper, 11pt]{exam}
\usepackage[left=3cm,bottom=2cm,top=3cm, right=2cm, headheight=2cm]{geometry}
% --- este trecho impede um conflito de nomes entre a classe exam e o pacote minted
\newcommand*{\renameenviron}[1]{%
	\expandafter\let\csname exam-#1\expandafter\endcsname
	\csname #1\endcsname
	\expandafter\let\csname endexam-#1\expandafter\endcsname
	\csname end#1\endcsname
	\expandafter\let\csname #1\endcsname\relax
	\expandafter\let\csname end#1\endcsname\relax
}
\renameenviron{framed}
\renameenviron{shaded}
\renameenviron{leftbar}
% ---
\usepackage[utf8]{inputenc}
\usepackage{graphicx}
\graphicspath{{../img/}}
% para posicionar equações à esquerda
\usepackage[fleqn]{mathtools} 
% syntax highlighting e formatação de código. impõe limitações no build deste documento. para mais detalhes veja a documentação do pacote.
\usepackage{minted}
% contador utilizado no cabeçalho "Página X de Y"
\usepackage{lastpage}
%\usepackage{showframe}
\usepackage{fancyvrb}

\firstpageheader{\includegraphics[height=1.7cm]{fapec}}{\textsc{Fundação Alagoana de Pesquisa, Educação e Cultura}\\\textsc{Faculdade de Tecnologia de Alagoas}\\}{\includegraphics[height=1.7cm]{fat}}

\firstpagefootrule
\firstpagefooter{}{}{\small Página \thepage \hspace{1pt} de \pageref{LastPage}}

\runningheader{\small Seleção para Monitoria em Programação}{}{}
\runningheadrule
\runningfooter{}{}{\small Página \thepage \hspace{1pt} de \pageref{LastPage}}
\runningfootrule

%Ajusta as palavras usadas na tabela de pontuação e questões
\vqword{Questão}
\vpword{Máximo}
\vsword{Pontuação}
\pointpoints{ponto}{pontos}

%-------------------------------------------------------------

\begin{document}
% -- cabeçalho
{
	\makebox[0.10\textwidth][l]{Curso:}
	\makebox[0.60\textwidth][l]{Análise e Desenvolvimento de Sistemas}
	\makebox[0.10\textwidth][l]{\enspace~Data:}
	\makebox[0.20\textwidth][l]{25/08/2017}
			
	\makebox[0.10\textwidth][l]{Prova:}
	\makebox[0.90\textwidth][l]{Monitoria em Introdução à Programação e Laboratório de Programação}
			
	\vspace{5mm}
	\makebox[0.10\textwidth][l]{Aluno:}
	\makebox[0.60\textwidth][l]{\hrulefill}
	\makebox[0.10\textwidth][l]{\enspace~Nota:}
	\makebox[0.15\textwidth][l]{\enspace\hrulefill}
			
	\vspace{5mm}
			
	% ---
			
	\textsc{Atenção candidato:} Todas as soluções devem ser escritas: (i) por \emph{completo}, isto é, nenhum trecho de código deve ser omitido ou abreviado, (ii) de forma legível e \emph{indentadas corretamente} e (iii) conter programas \emph{sintaticamente corretos}. A não observância destes critérios resultarão na redução da pontuação concedida.
			
	\begin{questions}
						
		\question[2\half] O somatório $0 + 1 + \ldots + n$ pode ser calculado utilizando a fórmula
		\begin{equation*} \label{sum}
			0 + 1 + \ldots + n = \frac{n \cdot (n + 1)}{2}
		\end{equation*}
						
		para $n > 0$. Escreva um programa que lê um número inteiro \texttt{n} e depois imprime o seu somatório utilizando a fórmula acima. Se o número lido for um valor menor que zero o programa não deve imprimir nada.
						
		\question[2\half] Escreva um programa que lê um número inteiro \texttt{n}, calcula e imprime seu fatorial. Considere que o número lido será sempre maior ou igual a zero. O fatorial de um número $n$ ($n > 0$), escrito $n!$, é:
		\[
			n! = 
			\begin{cases}
				n \cdot (n - 1) \cdot (n - 2) \cdot~\ldots~\cdot 2 \cdot 1, & \text{se } n > 1 \\
				1,                                                          & \text{se } n = 0 
			\end{cases}
		\]
						
		Exemplos:
		\begin{align*}
			5! & = 5 \cdot 4 \cdot 3 \cdot 2 \cdot 1 = 120 \\
			4! & = 4 \cdot 3 \cdot 2 \cdot 1 = 24          \\
			2! & = 2 \cdot 1 = 2                           \\
			0! & = 1                                       
		\end{align*}
						
		\question[2\half] Escreva um programa que, utilizando estruturas de repetição, preenche um array de capacidade 5 com números inteiros fornecidos pelo usuário. Em seguida o programa deve imprimir os números do array conforme o formato a seguir, também fazendo uso de estruturas de repetição. Por exemplo, se os números fornecidos forem \texttt{0}, \texttt{7}, \texttt{3}, \texttt{4} e \texttt{9}, nesta ordem, então o programa deve imprimir:
						
		\begin{Verbatim}[xleftmargin=\mathindent]
0␣␣␣9
␣7␣4␣
␣␣3␣␣
␣7␣4␣
0␣␣␣9
		\end{Verbatim}
						
		O caractere ␣ mostrado acima serve apenas para exibir a presença de um espaço em branco.
						
		\newpage
		\question[2\half] Escreva \emph{precisamente} o que imprime o código abaixo quando executado.
						
		\begin{minted}{java}
			public class WhatDoIDo {
				public static double f(int x, int y, int[] a) {
					double w = 0;
					for (; x <= y; y--) {
						System.out.println(w + " ");
						w += a[y];
					}
					System.out.println(w);
					return w / (y - x);
				}
				public static void main(String[] args) {
					int[] b = {1, 2, 3, 4, 5, 6, 7};
					b[b.length / 2] = b.length;
					f(1, b.length - 1, b);
				}
			}
		\end{minted}
						
	\end{questions}
			
\end{document}